\documentclass[11pt]{article}
\usepackage[top=1in,bottom=1in,left=1.1in,right=1.1in]{geometry}
\usepackage[T1]{fontenc}
\usepackage{lmodern}
\usepackage{hyperref}
\usepackage{setspace}

\hypersetup{colorlinks=true, urlcolor=black, linkcolor=black}
\pagestyle{empty}
\setstretch{1.08}
\setlength{\parskip}{0.8em}
\setlength{\parindent}{0pt}

\begin{document}

\textbf{\large Clovis SFEIR} \\
Nantes, France \\
+33\,6\,95\,83\,95\,34 \;$\cdot$\; \href{mailto:clovis.sfeir.cs@gmail.com}{clovis.sfeir.cs@gmail.com} \;$\cdot$\; \href{https://github.com/sfeirc}{github.com/sfeirc}

\vspace{1.2em}
D.\ E.\ Shaw \& Co. \\
Software Developer Internship -- New York, Summer 2027

\vspace{1.2em}
Dear D.\ E.\ Shaw Hiring Team,

I'm writing to apply for the Software Developer Intern position (New York, Summer 2027). D.\ E.\ Shaw's reputation for treating software engineering as a genuine research discipline -- building systems where correctness and performance both actually matter, not one traded off against the other -- is exactly the kind of environment I want to work in.

I'm an engineering student at IMT Atlantique (Nantes, France), on an apprenticeship track pairing my degree with an AI Engineer role at Infotel's Lab IA. Outside of coursework, I build systems I benchmark against real references rather than assumptions. Most relevant here: a C++20 limit-order-book matching engine implementing price-time priority with intrusive doubly-linked price-level lists for O(1) cancellation and pre-allocated object pooling to keep the hot path free of heap allocation -- it sustains 2.3M+ messages/second single-threaded with P50 latency of 0.6--1.0\,\textmu s, verified with a reproducible benchmark suite across three synthetic order-flow regimes, not a single cherry-picked run. A companion project models mean-reverting commodity spreads with an Ornstein-Uhlenbeck process, matched against a price-time-priority order book of my own -- I care about the same rigor on the modeling side as on the systems side, and verify one against the other.

I also placed 24th out of 124 teams (top 20\%, alongside entrants from Stanford, Princeton, and ETH Z\"urich) in the HRT/Partcl Macro Placement Challenge 2026, building a hybrid analytical VLSI placer combining force-directed global placement with timing-aware legalization.

What draws me to D.\ E.\ Shaw specifically is that it was founded as a computing company that happens to trade, not the reverse -- I'd rather spend a summer solving hard distributed-systems and performance problems where the domain itself is intellectually serious than doing generic backend work. I'd welcome the chance to bring that same build-it-and-verify-it discipline to your infrastructure.

Thank you for your consideration -- I'd be glad to walk through any of this work in more depth.

\vspace{1em}
Best regards, \\
Clovis Sfeir

\end{document}
